% 开头第一句：
% Directly using natural language (NL) to control arbitrary tables-to-tables transformations has long been a goal in the data community, as it can greatly improve the efficiency of data preparation (DP).

% 讲tables-to-tables transformations为什么难（从数据角度）？
% - 表格脏
% - 多表
% -

% - 步骤长
% - 自然语言的交互进一步加剧了任务难度

% ----------------

% Current Approaches and Their Limitations.
% - 总：受限于NL 的tables-to-tables transformations难度，现有方法都不能端到端的方式实现用自然语言控制 arbitrary tables-to-tables transformations 。
% - 一部分： focus on isolated sub-problems，例如数据清洗、数据merge...
% - 另一部分：提供更加明确的指导，例如提供目标表格的示例、提供固定的步骤指导（有workflow的指导）、预定义的固定操作

% 本质上，由于模型在长链任务上的建模能力弱，这些方法均是通过聚焦特定任务、或者提供预定义的规则，来降低tables transformations的建模难度。
% 落在任务复杂多样

% ----------------

% 幸运地是，大模型展现出了解决复杂任务的潜力。然而直接应用LLM仍然面临巨大的搜索空间，来源于两个方面：
% - 搜索空间大：自然语言输入和生成（生成代码）+ 链路长；整理可能搜索空间是：（词表大小）的（链路长度）次方
% - 缺少指导：没有全链路数据

% 搜索空间大+没有指导，模型无法学习

% ----------------

% 针对第一个问题，提出Tree-based Operation Orchestration的数据准备，将数据准备操作抽象出逻辑算子（类似于规范了一种交互范式）
% 针对第二个问题，我们提出数据合成

% ----------------

% Along this idea，我们提出DeepPrep，一个Unified End-to-End Data Preparation through Tree-based Operation Orchestration。

%  - 框架：
%     - 算子定义
%     - Tree-based Operation Orchestration
%  - 数据
%     - 前向
%     - 后向

\section{Introduction}
\label{sec:intro}

\begin{itemize}
    \item \textbf{Motivation and Problem Statement}
    \begin{itemize}
        \item The long-standing goal of using Natural Language (NL) to control arbitrary tables-to-tables (T2T) transformations for efficient Data Preparation (DP)~\cite{jin2025elt}.
        \item The inherent difficulty of this task, stemming from:
        \begin{itemize}
            \item Data complexity (\eg dirty data, multi-table integration).
            \item The added complexity of NL interaction.
            \item Consequently, the difficulty lies in mapping the complex, ambiguous NL instructions to the required multi-step, long-chain transformations.
        \end{itemize}
    \end{itemize}

    \item \textbf{Limitations of Current Approaches}
    \begin{itemize}
        \item Existing methods fail to achieve true end-to-end NL-\nlwhat arbitrary T2T transformations.
        \item Current solutions typically fall into two categories:
        \begin{itemize}
            \item Focusing on isolated sub-problems (\eg data cleaning, merging).
            \item Requiring extensive explicit guidance (\eg providing target examples~\cite{he2018transform}, fixed workflows~\cite{ge2025text}, or predefined, limited operations~\cite{DBLP:journals/pvldb/LiHYWC23}).
        \end{itemize}
        \item Fundamentally, these limitations stem from the inability of existing models to robustly \textit{compose and orchestrate} the multi-step, long-chain transformations required by diverse tasks, forcing these approaches to simplify the problem space through task isolation or explicit guidance.
    \end{itemize}

    \item \textbf{The LLM Opportunity and Core Challenge}
    \begin{itemize}
        \item Large Language Models (LLMs) show significant promise for solving complex, multi-step tasks.
        \item However, directly applying LLMs to T2T transformation faces a massive search space challenge, derived from:
        \begin{itemize}
            \item The vast combination of NL interpretation and code/operation generation.
            \item The exponential complexity of long transformation chains (Search Space $\approx (\text{Vocabulary Size})^{\text{Chain Length}}$).
        \end{itemize}
        \item This problem is compounded by a critical lack of full-chain supervision data for training.
    \end{itemize}

    \item \textbf{Our Solution: DeepPrep}
    \begin{itemize}
        \item We propose \textbf{DeepPrep}, a Unified End-to-End Data Preparation framework designed to overcome these challenges.
        \item DeepPrep introduces two core innovations:
        \begin{itemize}
            \item \textbf{Tree-Structured Compositional Policy Learning.} To tackle the exponential search space, we first define a set of \textit{logical operators} to standardize the T2T process. Building upon this, we propose a {Tree-based Operation Orchestration} mechanism, which significantly prunes the search space and enables the robust composition of complex transformation chains. We further enhance the model's ability to navigate this space through a novel {Curriculum Learning-based SFT+RL Training Strategy}.
            \item \textbf{Bi-Directional Full-Chain Supervision Synthesis.} To address the lack of supervision data, we introduce an innovative pipeline for generating full-chain training examples, featuring both {Forward Generation} and {Backward Generation}.
        \end{itemize}
    \end{itemize}

    \item \textbf{Contributions}
    \begin{itemize}
        \item We present the DeepPrep framework, detailing the definition of logical operators and the Tree-based Orchestration mechanism.
        \item We introduce our synthetic data generation pipeline.
        \item We demonstrate that DeepPrep achieves state-of-the-art performance in end-to-end NL-\nlwhat T2T transformation.
    \end{itemize}
\end{itemize}
